\documentclass[11pt,a4paper]{letter}
\usepackage[T1]{fontenc}
\usepackage{lmodern}
\usepackage[margin=2.4cm]{geometry}
\usepackage[hidelinks]{hyperref}
\usepackage{enumitem}
\signature{Yuan-chin Ivan Chang\\
Institute of Statistical Science, Academia Sinica\\
Taipei 11529, Taiwan\\
\href{mailto:ycchang@as.edu.tw}{ycchang@as.edu.tw}\\ORCID 0000-0002-4977-7721}
\address{Institute of Statistical Science\\Academia Sinica\\Taipei 11529, Taiwan}
\date{12 November 2026}

\begin{document}
\begin{letter}{Editor-in-Chief\\
\emph{Stochastic Environmental Research and Risk Assessment}}

\opening{Dear Editor,}

We submit the manuscript \emph{``Can the hidden state of a recurrent precipitation model monitor
drought onset? A pre-specified event-based evaluation of backward-coherence regularisation''} by
Hui-Mean Foo and Yuan-chin Ivan Chang for consideration as a Research Article in \emph{Stochastic
Environmental Research and Risk Assessment}.

Recurrent neural networks are increasingly used for hydrometeorological time series, and their hidden
states are an attractive source of monitoring statistics for environmental risks such as drought. We are
not aware, however, of a test of whether such learned statistics add anything to standard indices that
uses the statistical controls operational monitoring requires. Our paper asks this question for one well-defined
construction, a reverse-martingale-motivated penalty whose reconstruction defect is monitored by a CUSUM
procedure, and answers it with a design written down before any model was trained.

The study rests on daily gauge records for 1961--2024 from 133 stations in four contrasting climates
(central Texas, Germany, south-eastern Australia and Taiwan) and on 1{,}466 drought onsets in 2001--2024
defined from reanalysis root-zone soil moisture, a quantity that none of the detectors uses. Detectors are
compared at a matched alarm budget calibrated on separate years, scored at the event level, and compared
by paired cluster-bootstrap intervals that resample years and so respect the spatial dependence among
stations. An ablation computes the same defect without the penalty, and standard SPI and SPEI are fitted on
a 30-year reference period.

The main findings are as follows.
\begin{itemize}[leftmargin=1.5em,itemsep=2pt]
\item The penalty smooths the hidden-state path by 23.5\% while next-day forecast error stays equivalent,
within $\pm2\%$ by two one-sided tests, to the unregularised network and to a gated recurrent unit.
\item The defect carries a drought signal: it falls towards drought onset and alarms a median of 16 days
earlier than daily SPI-90 on events both detect. At the same alarm budget, however, it detects fewer onsets
(0.665 against 0.763; difference $-0.098$, 95\% CI $-0.140$ to $-0.056$). A 30-day SPI reproduces most of
its speed, and the defect adds no predictive information once SPI-30 is known.
\item A pre-registered extension that supplies temperature and evaporative demand improves the network's
forecasts but weakens the defect relative to SPEI (0.433 against 0.749 of onsets detected).
\item In neither form does the penalty improve monitoring.
\end{itemize}

We believe the paper suits the journal for two reasons. Its methodological content is stochastic and
statistical throughout: sequential change detection, calibration of monitoring procedures under
autocorrelation and seasonality, event-based verification, equivalence testing and dependence-aware
resampling, applied to a hydroclimatic risk. And its conclusions, including the negative ones, show how a
learned monitoring statistic can be benchmarked, against indices of matching time scale and against its
own ablation, before operational value is claimed. The framework can be applied to other learned
statistics, such as autoencoder reconstruction errors or forecast residuals of deep models.

We confirm that the manuscript is original, has not been published elsewhere and is not under
consideration by any other journal, and that both authors have approved the submission. The authors have no
competing interests. The work was partially supported by the National Science and Technology Council,
Taiwan (grant 115-2118-M-001-004-). All data are public, and the code that reproduces every number, table
and figure, together with the pre-specified protocol and its deviation log, is held in a repository that we
will make public upon acceptance and can share with the editor and reviewers on request. The use of AI tools
for code development and language editing is documented in the Methods, as the journal requires.

Thank you for considering our manuscript.

\closing{Sincerely yours,}

\end{letter}
\end{document}
